
Para poder hacer uso de las plenas capacidades de la computación \textit{real-time} para el estudio de los circuitos biohíbridos, no basta solo con alguno de los sistemas previamente mencionados, se necesita de herramientas específicas que hayan sido desarrolladas con el objetivo de interaccionar de manera determinista y en tiempo real con otros elementos. De las siguientes herramientas se han tenido en cuenta las que se usan comúnmente en la \textit{Escuela Politécnica Superior} de la \textit{UAM} para realizar experimentos con neuronas, o que hayan sido desarrollados por investigadores de esta. 
